% =============================================================================
%  lesson-note-template.tex — قالب مذكرة الدرس (المعلوماتية — التعليم الثانوي)
%  يُصرَّف بـ XeLaTeX حصراً. التصميم: أزرق/أسود/أبيض.
%
%  الاستعمال: انسخ هذا الملف إلى مجلد الدرس باسم lesson-note.tex
%             ثم عدّل القيم بين الأقواس المعقوفة (DEMO).
%  البناء (من مجلد الدرس):
%    xelatex -interaction=nonstopmode -halt-on-error -output-directory=build lesson-note.tex
% =============================================================================

\documentclass[11pt, a4paper]{article}

% ── المشترك (التمهيد + التصميم + الماكروهات) ────────────────────────────────
\input{"../../shared/common.tex"}

% ══════════════════════════════════════════════════════════════════════════════
\begin{document}

% ── الصفحة 1: الترويسة + البطاقة الفنية ─────────────────────────────────────
\officialheader

\cardtitle{مذكرة درس: <المورد المستهدف>}

{\gridsetup
\centering
\begin{tabularx}{0.85\textwidth}{@{}|c|X|@{}}
  \hline
  \cardrow{المستوى}{<المستوى، مثل: السنة الأولى ثانوي — جذع مشترك علوم وتكنولوجيا>}
  \cardrow{المادة}{<المادة، مثل: المعلوماتية>}
  \cardrow{التاريخ}{<التاريخ>}
  \cardrow{المدة}{<المدة، مثل: ساعة واحدة>}  \cardrow{المجال التعلمي}{<المجال التعلمي>}
  \cardrow{الوحدة التعلمية}{<الوحدة التعلمية>}
  \cardrow{المورد المستهدف}{<المورد المستهدف>}
  \cardrow{الكفاءة المستهدفة}{<الكفاءة المستهدفة>}
  \cardrow{الأهداف التعليمية}{<الأهداف التعليمية>}
  \cardrow{المكتسبات القبلية}{<المكتسبات القبلية>}
  \cardrow{الوسائل والدعامات}{<الوسائل والدعامات الديداكتيكية>}
\end{tabularx}\par
}\vspace{8pt}

% ── الصفحة 2: مراحل سير الوضعية التعلمية ────────────────────────────────────
\sectionbreak{مراحل سير الوضعية التعلمية}

\stageheader
\stagecontent
  {\badge{1}\\مرحلة الانطلاق}
  {<مهمات التعلم والهدف منها: التذكير بالمكتسبات القبلية وطرح الوضعية المشكلة>}
  {<نشاط الأستاذ والمتعلم>}
  {تقويم تشخيصي}
\stagecontent
  {\badge{2}\\بناء التعلمات}
  {<مهمات التعلم والهدف منها: اكتشاف المورد المعرفي الجديد وتطبيقه>}
  {<نشاط الأستاذ والمتعلم>}
  {تقويم تكويني}
\stagecontent
  {\badge{3}\\الحصيلة والتقويم}
  {<مهمات التعلم والهدف منها: صياغة الخلاصة وتثبيت المفاهيم>}
  {<نشاط الأستاذ والمتعلم>}
  {تقويم ختامي}

% ── الصفحة 3: الصياغة والتأسيس المعرفي ──────────────────────────────────────
\sectionbreak{الصياغة والتأسيس المعرفي}

\begin{rememberbox}[<العنوان المعرفي 1>]
  <التعريفات والمعارف…>
\end{rememberbox}

\begin{rememberbox}[<العنوان المعرفي 2>]
  <التعريفات والمعارف…>
\end{rememberbox}

% ── الصفحة 4: النشاط التطبيقي ───────────────────────────────────────────────
\sectionbreak{النشاط التطبيقي}

\begin{activitybox}[<التمرين الأول>]
  <نص التمرين…>
\end{activitybox}

\begin{activitybox}[<التمرين الثاني>]
  <نص التمرين…>
\end{activitybox}

\end{document}
